\answer{ex:06:1}{$e=(1-e^{-T_a/\tau_e})e^{-d/\tau_e}$. (가)~$0.110$ 대 $4.5\cdot10^{-5}$, 약 $2.5\cdot10^3$배. (나)~$\tau_e^*=T_a/\ln(1+T_a/d)=0.59\,\text{s}$.}
\answer{ex:06:2}{속도는 초당 $\nu_x\nu_y(A_+\tau_+-A_-\tau_-)=0.8A_+$, 한 시간에 약 $2900A_+$; $\tau_-=\tau_+/0.6=33\ms$.}
\answer{ex:06:3}{$\mu^2+s_1^2>s_2^2$이면 $\mathbf e_1$. 여기서는 $1.01>0.25$이므로, $\mathbf e_2$ 방향의 산포가 $25$배 큰데도 뉴런은 $\mathbf e_1$을 배운다. 이 규칙은 공분산행렬이 아니라 상관행렬의 주 벡터를 찾는다.}
\answer{ex:06:4}{전등과 주스 사이에서 $V=Re^{-(t_c+L-t)/\tau_\gamma}$이므로 $\dot V=V/\tau_\gamma$. 급증의 넓이는 $Re^{-L/\tau_\gamma}$: $0.61R$과 $0.78R$.}
\answer{ex:06:5}{신호 대 잡음비는 $1/(N+1)$, 시도 횟수는 $\propto N+1$: $5\cdot10^5$회 $\approx1$개월, $5\cdot10^{11}$회 $\approx8\cdot10^4$년.}
\answer{ex:06:6}{(가)~$k_2=1/\tau_d$, $k_1=1/\tau_d-1/\tau_\gamma$. (나)~거짓 오차는 $-(V/\tau_\gamma)\,\varepsilon/(1+\varepsilon)$, $\varepsilon=\tau_d/\tau_\gamma$이므로 $\tau_d\le0.105\,\text{s}$.}
\answer{ex:06:7}{$0.46$, $0.90$, $0.99$, $0.95$, 그리고 $d=3\,\text{s}$에서 $0.70$. 점들은 흔적 비율 $F=(1-e^{-T_{\text{cue}}/\tau_e})e^{-d/\tau_e}$의 한 곡선 위에 놓인다. $F\gtrsim0.1$이면 학습되고, $F<10^{-2}$이면 무작위 선택이다.}
\answer{ex:06:8}{$\eta^*=1/\bigl(2\sigma^2(N+2)\bigr)$, $N+2$회 시도. $N$개 중 $m$개만 흔들리면 $N(m+2)/m\ge N+2$회: 희소성은 도움이 되지 않는다.}
